\begingroup
\color{blue}
\begin{table}[H]
\centering
\caption{PV-plus-storage window-sensitivity diagnostic under the representative-day implementation.}
\scriptsize
\resizebox{\textwidth}{!}{%
\begin{tabular}{l r r r r L{0.38\linewidth}}
\toprule
Grid region & Central PV ratio & Earlier-window PV ratio & Later-window PV ratio & Max--min & Interpretation \\
\midrule
ERCOT & 0.00 & 0.00 & 0.00 & 0.00 & implementation diagnostic; no date-resolved shifted-calendar-week sensitivity was performed \\
MISO & 1.10 & 1.10 & 1.10 & 0.00 & implementation diagnostic; no date-resolved shifted-calendar-week sensitivity was performed \\
PJM & 0.50 & 0.50 & 0.50 & 0.00 & implementation diagnostic; no date-resolved shifted-calendar-week sensitivity was performed \\
SPP & 0.40 & 0.40 & 0.40 & 0.00 & implementation diagnostic; no date-resolved shifted-calendar-week sensitivity was performed \\
\bottomrule
\end{tabular}%
}
\tabnote{Values show the mid-2035, 4 h storage case. The current canonical implementation uses a representative JJA mean daily PV/load profile tiled over seven days; earlier/later labels therefore diagnose implementation limits and do not represent shifted calendar weeks. This diagnostic identifies the limitation of the representative-day implementation and motivates a date-resolved shifted-window sensitivity. Full storage-duration and scenario rows are retained in the source CSV.}
\end{table}
\endgroup
